\chapter*{Acknowledgements}

Certainly this book would have been possible without the support of
the National Science Foundation and of Five Colleges, Inc.  We
particularly want to thank Louise Raphael who, as the first director
of the calculus program at the National Science Foundation, had faith
in us and recognized the value of what had already been accomplished
at Hampshire College when we began our work.  Five College
Coordinators Conn Nugent and Lorna Peterson supported and encouraged
our efforts, and Five College treasurer and business manager Jean
Stabell has assisted us in countless ways throughout the project.

We are very grateful to the members of our Advisory Board: to Peter
Lax, for his faith in us and his early help in organizing and chairing
the Board; to Solomon Garfunkel, for his advice on politics and
publishing; to Barry Simon, for using our text and giving us his
thoughtful and imaginative suggestions for improving it; to Gilbert
Strang, for his support of a radical venture; to John Truxal, for his
detailed commentaries and insights into the world of engineering.

Among our colleagues, James Henle of Smith College deserves special
thanks.  Besides his many contributions to our discussions of
curriculum and pedagogy, he developed the computer programs that have
been so valuable for our teaching: Graph, Slinky, Superslinky, and
Tint.  Jeff Gelbard and Fred Henle ably extended Jim's programs to the
MacIntosh and to DOS Windows and X Windows.  All of this software is
available on anonymous ftp at emmy.smith.edu.  Mark Peterson, Robert
Weaver, and David Cox also developed software that has been used by
our students.

Several of our colleagues made substantial contributions to our
frequent editorial conferences and helped with the writing of early
drafts.  We offer thanks to David Cohen, Robert Currier and James
Henle at Smith; David Kelly at Hampshire; and Frank Wattenberg at the
University of Massachusetts.  Mary Beck, who is now at the University
of Virginia, gave heaps of encouragement and good advice as a
co-teacher of the earliest version of the course at Smith.  Anne
Kaufmann, an Ada Comstock Scholar at Smith, assisted us with extensive
editorial reviews from the student perspective.

Two of the most significant new contributions to this edition are the
appendix for graphing calculators and a complete set of solutions to
all the exercises.  From the time he first became aware of our
project, Benjamin Levy has been telling us how easy and natural it
would be to adapt our Basic programs for graphing calculators.  He has
always used them when he taught {\em Calculus in Context}, and he
created the appendix which contains translations of our programs for
most of the graphing calculators in common use today.  Lisa Hodsdon,
Diane Jamrog, and Marcia Lazo have worked long hours over an entire
summer to solve all the exercises and to prepare the results as
\LaTeX{} documents for inclusion in the Handbook for Instructors.  We
think both these contributions do much to make the course more useful
to a wider audience.

We appreciate the contributions of our colleagues who participated in
numerous debriefing sessions at semester's end and gave us comments on
the evolving text.  We thank George Cobb, Giuliana Davidoff, Alan
Durfee, Janice Gifford, Mark Peterson , Margaret Robinson and Robert
Weaver at Mount Holyoke; Michael Albertson, Ruth Haas, Mary Murphy,
Marjorie Senechal, Patricia Sipe, and Gerard Vinel at Smith.  We
learned, too, from the reactions of our colleagues in other
disciplines who participated in faculty workshops on Calculus in
Context.

We profited a great deal from the comments and reactions of early
users of the text.  We extend our thanks to Marian Barry at Aquinas
College, Peter Dolan and Mark Halsey at Bard College, Donald Goldberg
and his colleagues at Occidental College, Benjamin Levy at Beverly
High School, Joan Reinthaler at Sidwell Friends School, Keith Stroyan
at the University of Iowa, and Paul Zorn at St. Olaf College.  Later
users who have helped us are Judith Grabiner and Jim Hoste at Pitzer
College; Allen Killpatrick, Mary Scherer, and Janet Beery at the
University of Redlands; and Barry Simon at Caltech.

Dissemination grants from the NSF have funded regional workshops for
faculty planning to adopt Calculus in Context.  We are grateful to
Donald Goldberg, Marian Barry, Janet Beery, and to Henry Warchall of
the University of North Texas for coordinating workshops.

We owe a special debt to our students over the years, especially those
who assisted us in teaching, but also those who gave us the benefit of
their thoughtful reactions to the course and the text.  Seeing what
they were learning encouraged us at every step.
 
We continue to find it remarkable that our text is to be published the
way we want it, not softened or ground down under the pressure of
anonymous reviewers seeking a return to the mean.  All of this is due
to the bold and generous stance of W. H. Freeman.  Robert Biewen, its
president, understands---more than we could ever hope---what we are
trying to do, and he has given us his unstinting support.  Our
aquisitions editors, Jeremiah Lyons and Holly Hodder, have inspired us
with their passionate conviction that our book has something new and
valuable to offer science education.  Christine Hastings, our
production editor, has shown heroic patience and grace in shaping the
book itself against our often contrary views.  We thank them all.
